% !TEX root = ../main.tex
\chapter{The powdery, creased, smooth, round sphere}\label{ch:c07}
\begin{paracol}{2}

\begin{sources}
\begin{tabularx}{\linewidth}{@{}l l L@{}}
\toprule
\textbf{Video} & \textbf{Length} & \textbf{Content}\\
\midrule
\seg{7}{1}{Lecture 7: The powdery, creased, smooth, round sphere} & 1:53:06 & A single long worked example in place of new
theory: the sphere $S^2$ as a bare set (curves and functions, but no continuity); with the induced topology (continuity of a
curve, topological manifold, two charts, maximal atlas); with a smooth atlas (compatibility of the charts, smooth curves,
functions and maps, a tangent vector, a $(0,2)$ tensor field); and with a linear connection built from it (the metric, the
connection coefficients, homework: the curve is an autoparallel).\\
\bottomrule
\end{tabularx}

\smallskip
\textbf{Transcript} \file{materials/transcripts/C07-V01.txt} (German YouTube auto-captions).
\end{sources}

\begin{reading}[Where this lecture sits in the other notes]
\emph{Geometric Anatomy}: lectures 4 to 7 (topological spaces and the induced topology, topological manifolds, charts and
atlases, smooth atlases and maximal atlases; the same examples with the sphere). The tensor field and the connection of the
last two sections are used again in \cref{ch:c06} and in the lectures on connections of the other set.
\end{reading}

% ==================================================================
\section{Plan: stages of the sphere}\label{sec:sp-plan}

After consulting his conscience Schuller decided not to go on in the subject matter but to compute a 120 minute \emph{example}
in which everything done so far occurs, including all the computations a student had asked for in the last lecture (``how does
this work concretely?''). To his own astonishment he found, preparing it, that one only needs addition, multiplication and
division: ``whoever says there is too much maths here: it is really all very simple, since we now understand the structure''
\ts{7}{1}{0:00}.

The title names the successive stages of the sphere, the structures put on a set one after the other: first the sphere as a \emph{pure
set} (powdery, \emph{pulvrig}); with a topology it is a topological sphere (``creased'', \emph{knittrig}: it can be crumpled like
paper); with a smooth atlas, a smooth sphere (\emph{glatt}); and finally the round sphere (\emph{rund}), once we have added also a
parallel transport, i.e.\ a linear connection or covariant derivative. Equipping it step by step produces the round sphere
\ts{7}{1}{1:04}.

There is a pedagogical problem: when one says ``sphere'' one thinks already of the round one, one has only ever been shown round
spheres, and instead of building up structure, as a mathematician does, the physicist takes structure away. The only way would have been
a glass sphere to smash with a hammer (the idea came this morning in the shower, too late). Schuller e-mailed the man who prepares
the lecture demonstrations asking for a hollow glass, porcelain or clay sphere, ``at least from afar spherical'', to pulverise
in the lecture; the answer was ``there is a glass sphere in the collection, but not to be pulverised, since it is needed again and again
for the experiment on the density of air''. So it is a thought experiment \ts{7}{1}{2:05}.

% ==================================================================
\section{The powdery sphere: $S^2$ as a bare set}\label{sec:sp-powder}

\begin{definition}[New concepts: the sphere as a set]\label{def:sp-set}
As a \emph{bare set} (no further structure) define
\[
  S^2:=\bigl\{(m,n,p)\in\R^3\;\big|\;m^2+n^2+p^2=1\bigr\},
\]
the set of triples of real numbers with that property; e.g.\ $(1,0,0)$, $(0,-1,0)$, $(1/\sqrt2,\,1/\sqrt2,\,0)$ are elements
\ts{7}{1}{3:29}.
\end{definition}
One is probably so deformed by education that one thinks at once of a radius, of distances, \dots: but none of this is written there.
We only have a set with elements. Remember: the set $\{a,b\}$ and the set $\{b,a\}$ are equal, a set has no order, it is the
collection of its elements; one may shake it. The elements here are triples with special properties, but only a set
\ts{7}{1}{4:34}.

\begin{aside}[The hammer]
Draw the glass sphere with its childhood picture in mind (the sphere of our childhood, which is the round one). Now regard it
only as a pure set: let the hammer fall; what remains is a heap of dust. Fill the dust into a set container and hand
it over: ``this is $S^2$ as a bare set''. The points are still there, but they do not lie next to each other, none is
farther from another than any other: nothing; it is the pure set, one can shake it. This is what one has to do to destroy the
structure of the round sphere \ts{7}{1}{6:23}.
\end{aside}
\begin{coursenote}[What the picture of $\R^3$ smuggles in]
``What is the madness of saying that the $\R^3$ from which these pictures come is the space?'' (the origin of the $\R^3$ is
left aside). One learns at school to visualise $\R^3$ and with the visualisation many assumptions come in that one later thinks are
properties of $\R^3$, but are not. An element $(m,n,p)$ is an \emph{ordered triple of numbers}; the habit of drawing it as a point,
by agreeing secretly to go $m$ along the first axis, $n$ parallel to the second and $p$ parallel to the third, is a nice but
additional idea, not part of the triple. ``Do not always think that a triple is an element of three-dimensional space''
\ts{7}{1}{8:29}. A triple of numbers satisfying a further condition is just an element of a smaller set; always think ``dust, shake
it, still the same set'' \ts{7}{1}{10:14}.
\end{coursenote}

What can one do on the bare set? Already a lot.
\begin{example}[A curve and functions on the bare set]\label{ex:sp-curve}
A curve $\gamma\colon\R\to S^2$,
\[
  \gamma(\lambda):=(\cos\lambda,\;\sin\lambda,\;0).
\]
``But this is just a curve in $\R^3$?'' No: it lands in triples, and these triples satisfy the condition, so it is a curve in the
set $S^2$ \ts{7}{1}{7:45}. One can also define real functions on $S^2$ (``crazy stuff'') by giving a value to each admissible triple,
e.g.\ a formula involving $1+p$ in a denominator, defined everywhere except at the point $(0,0,-1)$ which one has to remove
\ts{7}{1}{10:35}. A map needs only sets \ts{7}{1}{11:17}.
\end{example}
Now the point. Paint the curve in blue chalk on the glass sphere of our childhood and smash it: the dust contains some grains painted
blue; shake the dust and hand it over: the points where the curve ran are scattered wildly over the powder. ``Is the curve
continuous?'' Here one cannot even talk about continuity; to follow the curve in the dust one would have to jump from grain to grain, with
nothing continuous about it \ts{7}{1}{11:37}. On a pure set continuity cannot be decided (we knew it already: continuity needs a topology),
much less differentiability, ``and straight line? God in heaven knows''. That is the powdery sphere \ts{7}{1}{12:59}.

% ==================================================================
\section{The creased sphere: a topology}\label{sec:sp-topological}

\subsection{Choosing a topology}

Now we make the topological sphere (here also called the creased sphere) \ts{7}{1}{13:20}. We must \emph{choose} some topology
and everything is open: the discrete topology (the power set, all subsets), or the chaotic one (only $\emptyset$ and $S^2$), the two
boring ones; we choose a third \ts{7}{1}{14:01}. ``There is nobody who forces you, but some topology we must choose.''

\begin{definition}[The induced (subset) topology on $S^2$]\label{def:sp-induced}
Let $\mathcal O_{\text{std}}$ be the standard topology of $\R^3$. As $S^2$ is a \emph{subset} of $\R^3$ we use the induced topology:
\[
  \mathcal O_{S^2}:=\bigl\{\,U\cap S^2\;\big|\;U\in\mathcal O_{\text{std}}\bigr\},
\]
the open sets of $S^2$ are those that can be written as the intersection of $S^2$ with an open set of $\R^3$ \ts{7}{1}{15:04}.
\end{definition}
A quick way to get a topology: $S^2$ inherits it from $\R^3$. Obviously another choice would give another topological sphere
\ts{7}{1}{16:14}. One of the reasons why this is clever is that it lets us decide quickly whether, for example, our curve $\gamma$ is
continuous \ts{7}{1}{16:34}. Unless said otherwise $\R^d$ is always equipped with the standard topology \ts{7}{1}{17:19}.

\subsection{Continuity of the curve}

\begin{theorem}[The curve $\gamma$ is continuous]\label{thm:sp-continuous}
With respect to $\mathcal O_{S^2}$ the curve $\gamma\colon\R\to S^2$ of \cref{ex:sp-curve} is continuous.
\end{theorem}
\begin{proof}
We look at the definition: we must show that for every open set the preimage under $\gamma$ is open \ts{7}{1}{17:40}. Let
$W\subseteq S^2$ be open. By definition of the induced topology there is an open $U\subseteq\R^3$ with $W=U\cap S^2$ \ts{7}{1}{18:42}.
The preimage is
\[
  \gamma^{-1}(W)=\{\lambda\in\R\mid\gamma(\lambda)\in U\cap S^2\}=\{\lambda\in\R\mid\gamma(\lambda)\in U\},
\]
because $\gamma$ lands in $S^2$ anyway; the intersection brings nothing new \ts{7}{1}{19:24}.
Now define an \emph{auxiliary curve} $\hat\gamma\colon\R\to\R^3$ (``a completely new curve, now in $\R^3$, allowed to lie anywhere''),
with $\hat\gamma(\lambda):=\gamma(\lambda)$, i.e.\ the same triples, now regarded as points of $\R^3$ (``a very subtle definition: mathematically
completely clear, but the point is that before we regarded the values in $S^2$ with its topology, and you can also regard them as a curve
in $\R^3$; if you wish you can forget the sphere''): $\hat\gamma(\lambda)=(\cos\lambda,\sin\lambda,0)$ \ts{7}{1}{21:00}. Then
$\gamma^{-1}(W)=\hat\gamma^{-1}(U)$. By school analysis (the $\varepsilon$-$\delta$ definition, equivalent to the topological one for the standard
topology) the curve $\hat\gamma\colon\R\to\R^3$ is continuous; so the preimage of the open set $U$ is open in $\R$, hence
$\gamma^{-1}(W)$ is open: $\gamma$ is continuous \ts{7}{1}{22:24}.
\end{proof}
Not hand-waving, then: we have seriously proved that the curve, as a curve in $S^2$ with the induced topology, is continuous. With another
topology it could have gone differently and the proof could not have been done this way. ``The nice thing about the standard topology:
if you imagine it embedded in something, what is lifted to $\R^3$ is also continuous in the subset topology; but one could do
otherwise; it is a pure matter of taste what one wants'' \ts{7}{1}{24:11}. In the same way one can show that the function $f$ above is
continuous on its domain: try it at home \ts{7}{1}{24:51}.
So this gave something: with a topology the sphere is not dust any longer, we can discuss whether a curve jumps \ts{7}{1}{25:13}.

\subsection{The topological manifold: two charts}

Next: this topological space is also a manifold. A topological manifold is a special case of a topological space with special
properties; so going from the topological sphere to a topological \emph{manifold} we put no new structure on it: either this topological space
\emph{is} one or it is not, and we have to check. One must show that around every point there is an open neighbourhood which can be mapped
continuously, with continuous inverse, onto chart images, and that the charts cover everything \ts{7}{1}{25:34}. So we
construct charts. Usually one takes longitude and latitude, but one is so used to these that one associates all sorts of things with
them; that is why the problem sheet used \emph{stereographic projection} (a plane at the north pole, one at the south pole)
\ts{7}{1}{26:58}. ``It looks as if everything depended on the sphere living in $\R^3$, but this construction is a crutch, a pure help for the
idea of constructing such charts''; to show it the definitions are simply written down and everything is verified from what is written
\ts{7}{1}{27:39}.

\begin{example}[The chart $x$]\label{ex:sp-chart-x}
Let $U:=S^2\setminus\{(0,0,-1)\}$ (a point of $S^2$ simply removed) and
\[
  x\colon U\to\R^2,\qquad (m,n,p)\mapsto\Bigl(\frac{2m}{1+p},\;\frac{2n}{1+p}\Bigr)
\]
(the first draft on the board had only one component; a student remarked that the second, $2n/(1+p)$, was missing; ``thank you, exactly''
\ts{7}{1}{29:44}). It is defined on all of $U$: at $p=-1$ the denominator would vanish, which is why the point $(0,0,-1)$ has to be removed and
the chart cannot be extended to it \ts{7}{1}{33:38}.
\end{example}
Is $x$ invertible? Write $(\alpha,\beta)$ for a pair in $\R^2$ and solve $\alpha=2m/(1+p)$, $\beta=2n/(1+p)$ for $(m,n,p)$, using
the third equation $m^2+n^2+p^2=1$: ``a task for small children: set equal, solve for three unknowns with three equations''
\ts{7}{1}{30:25}. The result is
\begin{equation}\label{eq:sp-x-inverse}
  x^{-1}(\alpha,\beta)=\Bigl(\frac{4\alpha}{\alpha^2+\beta^2+4},\;\frac{4\beta}{\alpha^2+\beta^2+4},\;\frac{4-\alpha^2-\beta^2}{\alpha^2+\beta^2+4}\Bigr)
\end{equation}
(``hopefully I computed this correctly this morning''; the auto-captions garble the third entry, which was re-derived here:
for $(\alpha,\beta)=(0,0)$ it must be the north pole $(0,0,1)$) \ts{7}{1}{31:07}.
Everything needs only multiplication and division; one solves a system of equations \ts{7}{1}{32:14}.
For this to be a chart, $x$ and its inverse must be continuous. $x$ is continuous if each component is, and the argument is the same as for the
curve, ``in the spirit of that one, with the preimage lying there and so on'': do it at home; the same for $x^{-1}$ \ts{7}{1}{32:57}.
We have covered almost the whole sphere with one chart, but only almost; so we need a second chart, \emph{because we want an atlas}.

\begin{example}[The chart $y$]\label{ex:sp-chart-y}
Let $V:=S^2\setminus\{(0,0,1)\}$ and
\[
  y\colon V\to\R^2,\qquad (m,n,p)\mapsto\Bigl(\frac{2m}{1-p},\;\frac{2n}{1-p}\Bigr),
\]
with inverse (writing $(\gamma,\delta)$ for a pair in $\R^2$ of the other chart, ``I could write $\alpha,\beta$ again; it makes no difference'')
\begin{equation}\label{eq:sp-y-inverse}
  y^{-1}(\gamma,\delta)=\Bigl(\frac{4\gamma}{\gamma^2+\delta^2+4},\;\frac{4\delta}{\gamma^2+\delta^2+4},\;\frac{\gamma^2+\delta^2-4}{\gamma^2+\delta^2+4}\Bigr)
\end{equation}
\ts{7}{1}{34:20}. (A student pointed out an error in the formula as first written; ``thank you; exactly, the inversion; the sign of the last
entry differs: a $+$ against a $-$''; the formula above is the corrected one \ts{7}{1}{35:23}.)
\end{example}
Again $y$ and $y^{-1}$ are continuous, as one sees by the same little auxiliary argument, ``since we chose this special topology''
\ts{7}{1}{36:05}.
The union $U\cup V$ is all of $S^2$. Around every point of $S^2$ we now have an open set ($U$ or $V$) which is mapped bijectively and
bicontinuously onto an open subset of $\R^2$. This shows: $(S^2,\mathcal O_{S^2})$ is a \emph{topological manifold} of dimension~$2$
(``the dimension of the manifold reads off from the charts''), the second stage of the sphere, the \textbf{creased sphere}
\ts{7}{1}{36:26}.

Why ``creased''? The childhood fantasy: the sphere is made of paper, crumple it; the curve is still continuous. You can make real
creases in the sphere, this changes nothing for the continuity of the curve. The round and the creased spheres are not distinguishable
topologically (``homeomorphic'': there is an invertible map between the two topological spaces, continuous in each direction). That is as it should be:
topology is about continuity and continuity is not attacked by creasing. (Pulverising it destroys the topology; ``going backwards is
much more fun'') \ts{7}{1}{37:32}.

\subsection{The maximal atlas}

A remark about the \textbf{maximal atlas} (``Max'' for short). The atlas $\mathcal A=\{(U,x),(V,y)\}$ just built is an atlas. Someone else
might come and say: ``I also put on the sphere the charts that I get from the same construction with the western and eastern pole, or with
any other point, I obtain many new charts and mine are also beautiful.'' You agree, as long as the continuity conditions in both directions
are satisfied. We saw already that if two charts are good for a topological manifold, the transition maps between them are automatically
continuous \ts{7}{1}{39:30}. So extend the atlas to the maximal one: it contains all charts $(W,z)$ for which the transition maps
\[
  z\circ x^{-1}\colon x(W\cap U)\to z(W\cap U),\qquad x\circ z^{-1}\colon z(W\cap U)\to x(W\cap U)
\]
are continuous, for \emph{all} charts $(U,x)$ of $\mathcal A$ (``one could also say: all charts compatible with all charts lying
there'') \ts{7}{1}{40:32}. The maximal atlas lets anybody who arrives crying ``my chart is also beautiful'' have it:
``yes, it belongs, as long as it is compatible with these'' \ts{7}{1}{41:56}.
What for? One does not really need it, but if one has a particular problem which would be much easier in a chart sitting at the
western pole, one may use it: it is in the maximal atlas. ``One always constructs a concrete atlas and then one may think the rest
along'' \ts{7}{1}{42:16}.

% ==================================================================
\section{The smooth sphere: a differentiable structure}\label{sec:sp-smooth}

\subsection{Selecting charts}

Not all transition maps in the maximal atlas are differentiable (there are continuous functions that are not differentiable) \ts{7}{1}{44:07}.
Hence the next stage, the differentiable (smooth) manifold \ts{7}{1}{44:28}. To make a topological manifold into a smooth one we
must add structure again, and this consists in \emph{selectively tearing charts out of the atlas} \ts{7}{1}{45:09}.
\begin{aside}[Information by removing]
One may encode a secret message by tearing out pages: ``Let the atlas have a million pages; number the letters of the alphabet from
$1$ to $26$; look at blocks of $26$ pages, and in each block tear out the page whose number is the letter''. So taking away gives
information: the tearing out of charts is structure \ts{7}{1}{45:30}.
\end{aside}

\subsection{Compatibility of the charts}

Go back to our innocent atlas $\mathcal A=\{(U,x),(V,y)\}$ and check whether it is a $C^\infty$ atlas (it could also be a $C^3$ atlas and so on):
we must check the differentiability class of the transition maps \ts{7}{1}{46:12}. First
$y\circ x^{-1}\colon x(U\cap V)\to y(U\cap V)$, and analogously, left to you, $x\circ y^{-1}$ \ts{7}{1}{46:54}.
Which mathematics decides this? The sets are subsets of $\R^2$ and the maps go from $\R^2$ to $\R^2$: standard analysis \ts{7}{1}{47:14}.
With \eqref{eq:sp-x-inverse} and the definition of $y$, the \emph{chart transition map} is
\begin{equation}\label{eq:sp-transition}
  (y\circ x^{-1})(\alpha,\beta)=\Bigl(\frac{4\alpha}{\alpha^2+\beta^2},\;\frac{4\beta}{\alpha^2+\beta^2}\Bigr),
\end{equation}
and $(x\circ y^{-1})(\gamma,\delta)$ has the same form with $(\gamma,\delta)$ \ts{7}{1}{47:56}.
(A terminology reminder: $x$ and $y$ are \emph{chart maps}; \eqref{eq:sp-transition} is a \emph{chart transition map}, a map $\R^2\to\R^2$ between
two chart images \ts{7}{1}{48:37}.)

Is it differentiable, and how often? Each component is rational, the only worry is where the denominator
$\alpha^2+\beta^2$ is zero. But $x(U\cap V)$ is the set $x(S^2\setminus\{\text{both poles}\})=\R^2\setminus\{(0,0)\}$, exactly the $\R^2$ without
the origin, so the map is smooth on its domain: it can be differentiated arbitrarily often \ts{7}{1}{49:39}. The inverse looks the same,
so the same holds. So $\mathcal A$ is a $C^\infty$ atlas \ts{7}{1}{50:19}.

One can again complete it to a maximal $C^\infty$ atlas $\mathcal A_{\max}^{\infty}$, adding all charts whose transition maps with the old ones are $C^\infty$.
It contains \emph{significantly fewer} charts than the maximal $C^0$ atlas, because the latter contains charts whose transition
maps are continuous but not differentiable, ``one can always build those'' \ts{7}{1}{51:01}.

\subsection{Smooth curves, functions, maps}

On a smooth manifold we can decide what we could not before: whether our curve $\gamma$, the one carried along since the very beginning, is
\emph{differentiable} with respect to this atlas; only if it is can we do all the beautiful things such as tangent vectors \ts{7}{1}{51:44}.
Recall the definition: choose a chart, if possible one in which the whole curve lies (otherwise one looks at several in turn: here there is
one). The question of smoothness reduces to that of the chart representation $x\circ\gamma\colon\R\to\R^2$.
\begin{example}[The curve is smooth]\label{ex:sp-smooth-curve}
The curve never leaves $U$ (never passes $(0,0,-1)$), so it can all be treated in the chart $x$. Composing, $\gamma(\lambda)=(\cos\lambda,\sin\lambda,0)$ has $p=0$, hence
\[
  (x\circ\gamma)(\lambda)=\bigl(2\cos\lambda,\;2\sin\lambda\bigr)\in\R^2 .
\]
This is a map $\R\to\R^2$ each of whose components is infinitely often differentiable (analysis book): $\gamma$ is $C^\infty$, a \emph{smooth curve}
\ts{7}{1}{55:35}. (To be rigorous one would need it for all charts of the atlas; the compatibility of the atlas takes care of it.)
``Everything is child's play once the definitions are known; you do not need the picture, our theory is complete; this does not mean you may not use intuition''
\ts{7}{1}{58:03}.
\end{example}
``Smooth is just the colloquial word: an object is called smooth if it is $C^\infty$, i.e.\ arbitrarily often continuously differentiable.'' Depending on the object
one decides in a different way: a curve is smooth if its chart representation is; a function $f\colon S^2\to\R$ is smooth if for all charts
the chart representation $f\circ x^{-1}$ is smooth; and a map between two smooth manifolds $\phi\colon M\to N$ is smooth if for all charts $(U,x)$ of $M$
and $(V,y)$ of $N$ the chart representation $y\circ\phi\circ x^{-1}\colon\R^d\to\R^{d'}$ is smooth (the two manifolds may have different
dimensions); using a second pair of charts $(\tilde x,\tilde y)$ one then checks it with the transition maps of the smooth atlases, so the answer
does not depend on the pair chosen \ts{7}{1}{58:26}. ``It is always the same principle: once understood, a nice exam question; you have no humour''
\ts{7}{1}{1:01:52}.

\subsection{A tangent vector}

\begin{example}[Tangent vector of $\gamma$]\label{ex:sp-tangent}
Compute the tangent vector $X_{\gamma,\gamma(0)}$ to $\gamma$ at $\lambda=0$ with the chart formula
\begin{align*}
  X_{\gamma,\gamma(0)}&=\dot\gamma^{\,m}_{(x)}(0)\,\Bigl(\frac{\partial}{\partial x^m}\Bigr)_{\gamma(0)},\\
  \dot\gamma^{\,m}_{(x)}&=(x^m\circ\gamma)' ,
\end{align*}
the $m$-th component of the chart representation differentiated at $\lambda=0$ \ts{7}{1}{1:04:11}. From the chart representation above,
\begin{align*}
  (x^1\circ\gamma)'(\lambda)&=-2\sin\lambda\;\to\;0,\ (x^2\circ\gamma)'(\lambda)&=2\cos\lambda\;\to\;2\qquad(\lambda=0),
\end{align*}
so $X_{\gamma,\gamma(0)}=2\,\bigl(\partial/\partial x^2\bigr)_{\gamma(0)}$: components $(0,2)$ in the basis induced by the chart $x$
\ts{7}{1}{1:08:06}. The point $\gamma(0)$ is the point $(1,0,0)$ \emph{of the manifold} $S^2$ (a triple, an element of the set), while its
image in the chart is $x(\gamma(0))=(2,0)\in\R^2$ \ts{7}{1}{1:08:27}.
\end{example}
\begin{coursenote}[Errors made and corrected in class]
On the board the derivative was first evaluated wrongly: the components came out as $(2,0)$ (i.e.\ $2\,\partial/\partial x^1$) because the
derivative was inserted without having differentiated $\cos$ and $\sin$ in the right order; after students' remarks (``this is the second basis vector'')
it was corrected: $m$ runs over the dimension of the manifold, $2$, and one is left with $2\,\partial/\partial x^2$. Moreover the base point had
been taken as $(0,0)$ instead of $(1,0,0)$: ``the point $(1,0,0)$ is a point of the manifold, an element of $S^2$''
\ts{7}{1}{1:06:59}. (The corrected values are used in the following.)
\end{coursenote}

% ==================================================================
\section{A $(0,2)$ tensor field}\label{sec:sp-tensor}

Now a ``useful'' $(0,2)$ tensor field \ts{7}{1}{1:08:48}. On a smooth manifold we have a tangent vector at every point, but a $(0,2)$ tensor field
is harder to imagine. Its component functions (always defined with respect to a chart) are smooth functions $g_{ij}\colon U\to\R$ with two indices,
$i,j=1,\dots,d$, here $d=2$: so in the chart $x$ the functions $g_{11},g_{12},g_{21},g_{22}$; to give a $(0,2)$ tensor field one must put four functions
on the table in each chart \ts{7}{1}{1:09:50}.

\begin{example}[A metric on $S^2$ given by brute force]\label{ex:sp-metric}
Functions on $S^2$ are best given as dependent on a point of the manifold, for which we can use the triple $(m,n,p)$ (the point lies in the
chart domain, a part of $S^2$). In the chart $x$ give
\begin{gather*}
  g_{(x)\,11}=g_{(x)\,22}=\frac{16}{(\alpha^2+\beta^2+4)^2},\\
  g_{(x)\,12}=g_{(x)\,21}=0,
\end{gather*}
where $(\alpha,\beta)=x(m,n,p)=\bigl(\tfrac{2m}{1+p},\tfrac{2n}{1+p}\bigr)$ are the chart coordinates of the point $(m,n,p)$ (in the lecture written out in $m,n,p$) \ts{7}{1}{1:12:21}. Is $g_{(x)11}$ smooth? Again just analysis, of
a function from $\R^3$ to $\R$: it is a quotient whose denominator is always positive, hence smooth on the domain \ts{7}{1}{1:13:23}.
That only covers the chart $x$ (it misses the point removed from $U$). On the other chart $y$ the same
functions are given with $1+p$ replaced by $1-p$:
\begin{gather*}
  g_{(y)\,11}=g_{(y)\,22}=\frac{16}{(\gamma^2+\delta^2+4)^2},\\ g_{(y)\,12}=g_{(y)\,21}=0,\quad(\gamma,\delta)=y(m,n,p)
\end{gather*}
\ts{7}{1}{1:14:45}.
\end{example}
\begin{coursenote}[Which pole?]
In the lecture the removed point of $U$ is first called ``the south pole'' and in the same breath the formulas for $y$ are said to cover the
other one; the consistent reading, used here, is: $x$ misses $(0,0,-1)$, $y$ misses $(0,0,1)$.
\end{coursenote}

``But wait: there was something about the transformation of tensor field components'' \ts{7}{1}{1:14:25}. If one just puts things on the table in two charts
one has not yet defined a tensor field: one must check that the components in the chart $y$ are obtained from those in chart $x$ by the
standard $(0,2)$ transformation law, otherwise the brute-force prescription would not define a tensor field.
This is the check one has to do \ts{7}{1}{1:16:23}.

\begin{theorem}[The prescription is a $(0,2)$ tensor field]\label{thm:sp-metric}
On the overlap $U\cap V$,
\[
  g_{(y)\,ij}\;=\;\frac{\partial x^m}{\partial y^i}\,\frac{\partial x^n}{\partial y^j}\;g_{(x)\,mn},
\]
the transformation law of a $(0,2)$ tensor field (\cref{thm:tf-tensor-trafo}), for $i,j\in\{1,2\}$.
\end{theorem}
\begin{proof}[Method, not all four components]
Both sides are evaluated at a point $(m,n,p)$ of the overlap, and we need $\partial x^m/\partial y^i$ at that point; by definition this is the $i$-th partial
derivative of the $m$-th component of $x\circ y^{-1}$, evaluated at $y(m,n,p)=(\gamma,\delta)$ \ts{7}{1}{1:17:48}. With $x\circ y^{-1}$ as in \eqref{eq:sp-transition} (first entry
$4\gamma/(\gamma^2+\delta^2)$), e.g.
\begin{align*}
  \frac{\partial x^1}{\partial y^1}&=\frac{4(\gamma^2+\delta^2)-8\gamma^2}{(\gamma^2+\delta^2)^2}\\
  &=\frac{4\delta^2-4\gamma^2}{(\gamma^2+\delta^2)^2}
\end{align*}
(product rule) \ts{7}{1}{1:20:35}. One computes in the same way $\partial x^1/\partial y^2$, $\partial x^2/\partial y^1$, $\partial x^2/\partial y^2$, and for each of the
four pairs $(i,j)$ one sums over $m,n$ on the right (four summands): ``you need four computations, and if everything has been done right, the left side, the given $g_{(y)\,ij}$,
comes out''. If it does come out for all four, we have defined a tensor field; if not, the brute-force prescription would not have been a $(0,2)$ tensor
field \ts{7}{1}{1:21:57}.
\end{proof}

\subsection{What a $(0,2)$ tensor field is good for}

``Who wants a $(0,2)$ tensor field?'' Everybody does, all the time. Stand at the north pole of the smooth sphere; a car drives by, and at the same point
another with a different velocity. (Velocities are not arrows, they are directional derivatives, but let us imagine for a moment.) One would like to say that the
velocity of one is perpendicular to the other's; but to say ``perpendicular'' one needs a scalar product, an inner product: it takes two vectors and gives a
number. Or one takes one vector and inserts it twice: that gives the square of the speed \ts{7}{1}{1:24:01}. And this should not depend on the chart: ``how should another
chart change the objective speed of the car? Objectively not at all''. For this one needs a $(0,2)$ tensor field \ts{7}{1}{1:25:23}.

\begin{definition}[New concepts: the number $g(X,Y)$, independence of the chart]\label{def:sp-gXY}
Let $g_{ij}$ be the components of a $(0,2)$ tensor field in a chart, and $X^i$, $Y^j$ the components of two vector fields in the same chart. Then
\[
  g_{ij}\,X^i\,Y^j
\]
(four summands for $d=2$) is a chart-independent real number (``chart-independent real number''; in the lecture the definition is stated with the symbol
$g$ on its own) \ts{7}{1}{1:26:10}.
\end{definition}
\begin{proof}
Everything is chart-dependent on the face of it (``it says everywhere `with respect to the chart'''), but look at a change to another chart: the components
transform as $g_{(y)ab}=\dfrac{\partial x^i}{\partial y^a}\dfrac{\partial x^j}{\partial y^b}g_{(x)ij}$ and $X^a_{(y)}=\dfrac{\partial y^a}{\partial x^m}X^m_{(x)}$,
$Y^b_{(y)}=\dfrac{\partial y^b}{\partial x^n}Y^n_{(x)}$. Then in $g_{(y)ab}X^a_{(y)}Y^b_{(y)}$ each Jacobian pair $\dfrac{\partial x^i}{\partial y^a}\dfrac{\partial y^a}{\partial x^m}$ gives a Kronecker
delta $\delta^i{}_m$ (\cref{thm:tf-inverse}), and one gets $g_{(x)mn}X^m_{(x)}Y^n_{(x)}$: the same expression in the other system \ts{7}{1}{1:27:34}.
\end{proof}
So the four numbers at each point define \emph{an inner product}: in a linear vector space with a basis one has an inner product, and now on our smooth manifold
there is a tangent space at each point, and in each tangent space an inner product given by $g$ \ts{7}{1}{1:29:05}.
Confusion arises from the habit of writing the inner product of vectors as $v^1w^1+\dots+v^dw^d$: this definition has a serious problem; it is fine in a
vector space and for a fixed basis, but it corresponds to the components $(\begin{smallmatrix}1&0\\0&1\end{smallmatrix})$ of $g$, and under a change of coordinates this
beautiful unit matrix becomes something wild. ``One tries to tell you that there is the standard inner product, the identity matrix \dots'': one cannot simply \emph{give numbers}, the
result should be chart-independent, so one must give a $(0,2)$ tensor field, which transforms exactly opposite to the vector components, so that everything is invariant
\ts{7}{1}{1:29:49}. Then, putting $Y=X$, one defines the square of the length of a vector with respect to this $(0,2)$ tensor field \ts{7}{1}{1:31:52}.

\begin{aside}[The metric and the length of curves]
This is only an example of the use of a $(0,2)$ tensor field, but it is the most important one: in differential geometry it is called a \textbf{metric}. With it one can
also measure the lengths of curves: one only said the curve is differentiable, but what is its total length? One has to integrate $\sqrt{g(\dot\gamma,\dot\gamma)}$ along it
(the lecture only says: integral over $g$ with the tangent vector of the curve plugged in). All this later \ts{7}{1}{1:32:35}.
\end{aside}

A student asks (looking sad) whether there is no such tensor field of this kind that comes for free. Answer, ``a very good question'': no. On a smooth manifold the length, e.g.\ of a
tangent vector, is not defined, just as in linear algebra the norm of a vector is not defined as long as no inner product is given. Length and angle need an inner
product. Example: ``the polynomials up to a certain degree form a vector space; what is the length of $x^2$? What should that mean? One must first give an inner product
on that space''. The same on the smooth manifold: the smooth sphere you picture as round, but the smooth manifold is also the sphere pulled long like a banana, as long as it has no
creases; as one can stretch it like rubber without making creases, the lengths would change; so there is no length as long as no $(0,2)$ tensor field is given. ``You imagine too much,
when you say `but there is a concrete distance, and the vector must have this length'. A length is always with respect to a $(0,2)$ tensor field, and one requires it to be \emph{symmetric}:
$g_{ij}=g_{ji}$, as in our example.'' Again: you have to give this structure, otherwise you cannot talk about it \ts{7}{1}{1:33:42}.

\subsection{The speed of the curve}

Compute the speed of the curve $\gamma$ at $\gamma(0)$ with respect to the $g$ just written down \ts{7}{1}{1:36:11}. We need the components of the tangent vector, $(X^1,X^2)=(0,2)$, and
$g$ at the point $(1,0,0)$ (not at $(0,0)$ nor $(0,0,1)$: the point of the manifold, $x=(2,0)$). Then
\[
  g\bigl(X,X\bigr)=g_{11}X^1X^1+g_{12}X^1X^2+g_{21}X^2X^1+g_{22}X^2X^2=g_{22}\,(X^2)^2,
\]
since $X^1=0$ and $g_{12}=g_{21}=0$. At $(m,n,p)=(1,0,0)$, i.e.\ $(\alpha,\beta)=(2,0)$:
\[
  g_{22}=\frac{16}{(2^2+0+4)^2}=\frac{16}{64}=\frac14,\qquad g(X,X)=\frac14\cdot 2\cdot2=1 .
\]
\ts{7}{1}{1:37:33}. So the squared speed is $1$: with respect to this strange $(0,2)$ tensor field the speed of the curve at that point is $1$ \ts{7}{1}{1:41:46}.
\begin{coursenote}[Errors made and corrected in class]
First the point was inserted wrongly ($m=n=0$, i.e.\ the north pole instead of $(1,0,0)$), and ``$16/4^2\cdot4=4$'' came out where $1$ was expected; ``this is wrong, excuse me,
I said earlier already the point $\gamma(0)$ is $(1,0,0)$''. With the right point the value is $1$ (``I am happy it is right'') \ts{7}{1}{1:40:23}.
The denominator in the lecture, $2^2+4=8$, squared gives $64$; $16/64\cdot4=1$ as above. Marked here since the transcript has the intermediate numbers garbled [?].
\end{coursenote}
What was done tacitly: this $g$ is the \textbf{pull-back} of the standard $(0,2)$ tensor field of $\R^3$ (the Euclidean inner product) to the sphere, expressed in the chart: it is the metric one gets by
laying a thread in $\R^3$ over a sphere and asking how far it is from there to there; not explained further here (Schuller: ``I secretly did what I did not tell you how to do: I inherited the
metric from $\R^3$''). The lesson: how long a vector is depends on the chosen $(0,2)$ tensor field \ts{7}{1}{1:42:07}.

% ==================================================================
\section{The round sphere: a linear connection from the metric}\label{sec:sp-round}

Finally: we showed how to give a $(0,2)$ tensor field brutally. One may also construct one from fields already given \ts{7}{1}{1:42:48}; we do it now for the next interesting object, the \textbf{linear connection} of \cref{ch:c06}.
Also there the components must be obtained in some way: in every chart $\Gamma^a{}_{bc}$ at every point, which must then transform correctly with the extra inhomogeneous term (\cref{def:cf-connection}).
In $d=2$ there are $2\cdot2\cdot2=8$ functions to give and to check; instead we do something clever: from the $(0,2)$ tensor field, our length measure, we \emph{construct} a linear connection
\ts{7}{1}{1:43:31}.

\begin{definition}[New concepts: the connection built from a $(0,2)$ tensor field]\label{def:sp-christoffel}
Let $g_{ab}$ be a symmetric, invertible $(0,2)$ tensor field, and $(g^{-1})^{ab}$ the $(2,0)$ tensor field with $(g^{-1})^{am}g_{mb}=\delta^a{}_b$ (on the sheet: ``the inverse matrix'', a tensor field
by the standard check; \ts{7}{1}{1:45:41}). Define in every chart
\[
  \Gamma^a{}_{bc}\;:=\;\tfrac12\,(g^{-1})^{am}\Bigl(\frac{\partial}{\partial x^b}g_{mc}+\frac{\partial}{\partial x^c}g_{bm}-\frac{\partial}{\partial x^m}g_{bc}\Bigr)
\]
\ts{7}{1}{1:46:22}. (The factor $\tfrac12$ is not audible in the auto-captions; it is the usual one and is needed for the transformation law, which gives the right inhomogeneous term [?].)
That these $\Gamma^a{}_{bc}$ \emph{do} transform as a linear connection (the tensorial part plus the inhomogeneous term of \eqref{eq:cf-gamma-trafo}) is a computation on problem sheet~5 (exercise~3, ``these are
connection components''): ``you know the transformation behaviour of the $g$'s and of the derivatives of anything; insert it all, and what remains beyond the tensor behaviour is the extra term''
\ts{7}{1}{1:47:08}.
\end{definition}
For our $g$ one obtains by inserting (``you can really do it at home, you only insert''), with $r^2:=\alpha^2+\beta^2$ the chart quantities in the chart $x$ \ts{7}{1}{1:47:53},
\begin{align*}
  \Gamma^1{}_{11}&=-\Gamma^1{}_{22}=\Gamma^2{}_{12}=\Gamma^2{}_{21}=-\frac{2\alpha}{r^2+4}
    ,\\
  \Gamma^2{}_{22}&=-\Gamma^2{}_{11}=\Gamma^1{}_{12}=\Gamma^1{}_{21}=-\frac{2\beta}{r^2+4}
\end{align*}
(in the lecture, the three coefficients $\Gamma^1{}_{11}$, $-\Gamma^1{}_{22}$, $\Gamma^2{}_{12}$ are given and a second quadruple ``$\Gamma^1{}_{12}$, $-\Gamma^2{}_{11}$, $\Gamma^2{}_{22}$'', each with the factor
$2\,x/\bigl(r^2+4\bigr)$ in terms of $m,n,p$; the signs and the $\beta$ group were not audible and are \emph{derived here} from the formula for $g=f\,\delta$ with
$f=16/(r^2+4)^2$, hence $\Gamma^a{}_{bc}=\tfrac12\bigl(\delta^a{}_b\partial_c\ln f+\delta^a{}_c\partial_b\ln f-\delta_{bc}\partial_a\ln f\bigr)$ [?]). The symmetry
$\Gamma^a{}_{bc}=\Gamma^a{}_{cb}$ means there are not eight independent functions, ``only about four types'' \ts{7}{1}{1:49:03}.
``You wanted number examples; I did not give them to you: so we construct a linear connection and check that it is one'' \ts{7}{1}{1:50:04}.

\begin{intuition}[Metric versus connection]
A metric measures: if the length of a tangent vector can be computed, the length of a whole curve follows by integration (later). A linear connection says how to ``follow one's nose'':
which are the straightest curves. Given a metric one can pick out the shortest curves; given a connection one picks out the straightest. This particular construction of $\Gamma$ from $g$ is special:
it makes the straightest curves also the shortest ones. ``You do not have to do this; but if you do, the shortest curves are the straightest ones; one may want that; I only tell you the background''
\ts{7}{1}{1:50:25}.
\end{intuition}

\paragraph{The only homework of the day.} Show that the curve $\gamma(\lambda)=(\cos\lambda,\sin\lambda,0)$ carried through the lecture from the very beginning is an \emph{autoparallel} with respect to
the linear connection so constructed (\cref{def:cf-autoparallel-curve}): write down the autoparallel equation $\ddot q^a+\Gamma^a{}_{bc}\dot q^b\dot q^c=0$ for $q=x\circ\gamma$, insert everything, ``compute, compute,
compute, half a page, and suddenly $0$ stands there; ``aha, it is an autoparallel''; and this can only be done on a sphere with a linear connection: this sphere has now a \emph{shape}, and it is the
sphere one knows intuitively'' \ts{7}{1}{1:51:47}. ``Until Monday'' \ts{7}{1}{1:52:49}.
\begin{coursenote}[A check of the homework]
Added, not said in class: with $q(\lambda)=(2\cos\lambda,2\sin\lambda)$ ($r^2=4$) and the $\Gamma$ above one finds for the first component
$\ddot q^1+\Gamma^1{}_{11}(\dot q^1)^2+2\Gamma^1{}_{12}\dot q^1\dot q^2+\Gamma^1{}_{22}(\dot q^2)^2
=-2\cos\lambda-2\cos\lambda\sin^2\lambda+4\sin^2\lambda\cos\lambda+2\cos^3\lambda=0$, and analogously for the second one (the equator is a great circle).
\end{coursenote}

% ==================================================================
\section*{Key concepts}
\begin{keyconcepts}
\textbf{Retrieval questions.}
\begin{enumerate}
  \item What can and what cannot be decided on $S^2$ as a bare set? Why can one not talk about continuity?
  \item Which topology did we put on $S^2$, and how does one prove that the curve $\gamma$ is continuous in it?
  \item Write the two charts $x,y$, their inverses and the transition map. Why is the atlas $\{(U,x),(V,y)\}$ a $C^\infty$ atlas?
  \item What is the maximal atlas, and why do the maximal $C^0$ and $C^\infty$ atlases differ?
  \item How does one decide smoothness of a curve, of a function, and of a map between manifolds?
  \item Compute the tangent vector to $\gamma$ at $\lambda=0$ in the chart $x$.
  \item Why does $g_{ij}X^iY^j$ not depend on the chart, and why does one need a $(0,2)$ tensor field to define lengths?
  \item How is a linear connection built from a metric, and why is the curve $\gamma$ an autoparallel?
\end{enumerate}
\textbf{Concepts.}
\begin{itemize}
  \item The stages of the sphere: set, topology (\cref{def:sp-induced,thm:sp-continuous}), topological and smooth manifold (charts \cref{ex:sp-chart-x,ex:sp-chart-y}; transition map \eqref{eq:sp-transition}).
  \item Smooth curve and tangent vector (\cref{ex:sp-smooth-curve,ex:sp-tangent}).
  \item $(0,2)$ tensor field as metric (\cref{ex:sp-metric,thm:sp-metric,def:sp-gXY}).
  \item Linear connection from the metric (\cref{def:sp-christoffel}); autoparallel curve (\cref{ch:c06}).
\end{itemize}
\end{keyconcepts}

\flushside
\end{paracol}
